\documentclass[12pt,a4paper]{article}

\usepackage[utf8]{inputenc}
\usepackage[T1]{fontenc}
\usepackage{amsmath,amssymb,amsthm}
\usepackage{algorithm}
\usepackage{algpseudocode}
\usepackage{booktabs}
\usepackage{array}
\usepackage{multirow}
\usepackage{graphicx}
\usepackage{xcolor}
\usepackage{hyperref}
\usepackage{geometry}
\usepackage{enumitem}
\usepackage{titlesec}
\usepackage{fancyhdr}
\usepackage{cite}
\usepackage{listings}

\geometry{margin=2.5cm}

\hypersetup{
  colorlinks=true,
  linkcolor=blue!70!black,
  citecolor=green!60!black,
  urlcolor=blue!70!black
}

\lstset{
  basicstyle=\ttfamily\small,
  breaklines=true,
  frame=single,
  backgroundcolor=\color{gray!8},
  keywordstyle=\color{blue},
  commentstyle=\color{green!60!black},
  stringstyle=\color{red!70!black}
}

\pagestyle{fancy}
\fancyhf{}
\rhead{S$^3$-Vote}
\lhead{Secure, Scalable, Sovereign Voting}
\rfoot{Page \thepage}

\theoremstyle{definition}
\newtheorem{definition}{Definition}[section]
\newtheorem{theorem}{Theorem}[section]
\newtheorem{property}{Property}[section]

% ─────────────────────────────────────────────────────────────────────────────
\title{
  \vspace{-1cm}
  \textbf{\Large S$^3$-Vote: A Reputation-Weighted, Blockchain-Anchored\\
  Decentralised Voting System with Adaptive Reconstruction Threshold}\\[0.6em]
  \large A Threshold Cryptography Framework with Publicly Verifiable Secret Sharing
  and Dynamic Trust Management
}
\author{Milad Seifi}
\date{\today}
% ─────────────────────────────────────────────────────────────────────────────

\begin{document}

\maketitle
\thispagestyle{fancy}

\begin{abstract}
Secure electronic voting requires confidentiality of individual ballots,
verifiability of the final tally, and resilience against colluding or
misbehaving participants — all without relying on a single trusted authority.
We present \textbf{S$^3$-Vote} (Secure, Scalable, Sovereign Voting), a
decentralised voting system that combines ElGamal threshold encryption,
Shamir Secret Sharing, and Publicly Verifiable Secret Sharing (PVSS) with a
novel \emph{reputation-weighted adaptive reconstruction threshold}.
Each decryption authority accumulates a persistent reputation score derived
from their verifiable behaviour across elections.
Reputation is converted to a weight, and the effective quorum required for
reconstruction adapts automatically: trusted groups need fewer participants;
distrusted or heterogeneous groups require more.
A two-gate reconstruction policy — enforcing both a cryptographic share
count and a reputation weight threshold simultaneously — closes attack
vectors that single-threshold schemes leave open.
All events, encrypted votes, PVSS commitments, share submissions, cheat
evidence, and reputation updates are anchored on a hash-chained blockchain,
providing a tamper-evident, non-repudiable audit trail.
Our scheme improves upon Shamir's foundational work and Nojoumian's
social secret sharing by eliminating the trusted dealer, providing
publicly verifiable cheat detection, and introducing a self-governing
trust mechanism that adapts dynamically across election rounds.
\end{abstract}

\tableofcontents
\newpage

% =============================================================================
\section{Introduction}
% =============================================================================

Democratic and organisational decision-making increasingly relies on
electronic voting systems. Yet most deployed systems route all trust
through a single entity — a government authority, a corporate IT department,
or a cloud provider. This creates two fundamental vulnerabilities: a
\emph{single point of trust} (the operator can manipulate results
undetected) and a \emph{single point of failure} (if the operator is
compromised, the election is ruined).

Threshold cryptography distributes trust across $N$ decryption authorities:
the private key required to decrypt ballots is split into $N$ shares, and
any coalition of at least $t$ shares can reconstruct it, while fewer than
$t$ shares reveal nothing. This removes the single point of trust.
However, classical threshold schemes treat all shares as equal.
A malicious authority submitting a corrupted share causes silent
reconstruction failure, and the system has no mechanism to identify
or penalise the attacker, nor to adapt future requirements based on
observed behaviour.

S$^3$-Vote addresses these limitations through five integrated contributions:

\begin{enumerate}[label=\textbf{\arabic*.}]
  \item \textbf{PVSS-based share verification} — each submitted share is
    publicly verifiable against an on-chain commitment, making cheating
    immediately detectable and undeniable.
  \item \textbf{Persistent reputation scoring} — authority behaviour
    is scored across Participation, Contribution, and Honesty dimensions
    and accumulated across elections with exponential smoothing.
  \item \textbf{Reputation weights} — scores are normalised into per-authority
    weights, so high-reputation authorities carry proportionally more
    influence over reconstruction.
  \item \textbf{Adaptive two-gate quorum} — reconstruction requires satisfying
    both a cryptographic share count gate \emph{and} a reputation weight
    gate simultaneously; the weight gate adapts dynamically based on
    collective trust and evenness.
  \item \textbf{Blockchain audit trail} — all events are permanently
    recorded on a hash-chained ledger, providing non-repudiable evidence
    of every action.
\end{enumerate}

% =============================================================================
\section{Background and Related Work}
% =============================================================================

\subsection{ElGamal Encryption}

ElGamal encryption \cite{elgamal1985} operates over a cyclic group
$\mathbb{Z}_p^*$ with generator $g$ and prime $p$.
A key pair $(y, x)$ is generated where $y = g^x \bmod p$ is the public key
and $x$ is the private key.
To encrypt a message $m$ (vote option, encoded as a group element):

\begin{align}
  C_1 &= g^k \bmod p \\
  C_2 &= m \cdot y^k \bmod p
\end{align}

where $k$ is a random ephemeral value.
Decryption recovers $m = C_2 \cdot C_1^{-x} \bmod p$.
The scheme is semantically secure under the Decisional
Diffie-Hellman (DDH) assumption.

In S$^3$-Vote, the private key $x$ is never held by any single party;
it is split via Shamir Secret Sharing before the election begins and
reconstructed only after all votes are cast.

\subsection{Shamir Secret Sharing}

Shamir's $(t,n)$-threshold scheme \cite{shamir1979} splits a secret
$s \in \mathbb{Z}_p$ by constructing a random polynomial of degree $t-1$:
\begin{equation}
  f(z) = s + a_1 z + a_2 z^2 + \cdots + a_{t-1} z^{t-1} \pmod{p}
\end{equation}
Authority $i$ receives share $s_i = f(i)$.
Any $t$ shares uniquely determine $f$ (via Lagrange interpolation) and
hence reveal $s$; any $t-1$ shares are information-theoretically independent
of $s$.

The classical scheme treats all shares as interchangeable: any $t$-subset
suffices. S$^3$-Vote augments this with reputation weights so that the
effective contribution of each share depends on the authority's trust level.

\subsection{Publicly Verifiable Secret Sharing (PVSS)}

In standard secret sharing, a participant who submits a corrupted share
causes reconstruction to fail with no publicly verifiable evidence of
wrongdoing. PVSS \cite{schoenmakers1999} solves this by publishing
\emph{commitments} alongside shares such that anyone can verify correctness
without learning the secret.

In S$^3$-Vote, at election creation, each share $s_i$ generates a commitment:
\begin{equation}
  C_i = g^{H(s_i)} \bmod p
\end{equation}
where $H$ is a cryptographic hash function.
These commitments are written to the blockchain before reconstruction begins.
When authority $i$ submits $\hat{s}_i$ during reconstruction, the verifier
checks $g^{H(\hat{s}_i)} \equiv C_i \pmod{p}$.
A mismatch is on-chain proof of submission of an invalid share.

\subsection{Reputation in Secret Sharing: Nojoumian's Scheme}

Nojoumian \textit{et al.}\ \cite{nojoumian2012,nojoumian2018} introduced the
concept of \emph{social secret sharing}, where participant reputation
influences the threshold. Their key contributions were recognising that
not all participants should be treated equally and that past behaviour
should inform future requirements.

However, their scheme has several limitations that S$^3$-Vote addresses:
\begin{itemize}
  \item Requires a \emph{trusted dealer} to assign and update reputation scores
  \item Reputation updates are manual or social (participants evaluate peers),
    introducing subjectivity and gamesmanship
  \item No mechanism for public proof of cheating
  \item Single reconstruction gate (no two-gate policy)
  \item No protection against reputation concentration
  \item No blockchain for tamper-evident history
  \item Designed as a general secret sharing primitive, not a complete
    voting system
\end{itemize}

% =============================================================================
\section{System Architecture}
% =============================================================================

\subsection{Participants}

\begin{itemize}
  \item \textbf{Election Administrator} — creates the election, defines
    authority set, options, and voter count.
  \item \textbf{Voters} — cast encrypted ballots; their identities are
    independent of the decryption process.
  \item \textbf{Decryption Authorities} ($A_1, \dots, A_N$) — each holds
    one Shamir share of the private key; collectively responsible for
    reconstruction after voting closes.
  \item \textbf{Blockchain} — append-only tamper-evident ledger; records
    all events and acts as the public bulletin board.
\end{itemize}

\subsection{Election Lifecycle}

The full protocol proceeds through five phases:

\begin{enumerate}
  \item \textbf{Setup} — key generation, secret splitting, PVSS commitment
    publication, threshold computation.
  \item \textbf{Voting} — voters submit ElGamal-encrypted ballots; ciphertexts
    recorded on-chain.
  \item \textbf{Reconstruction} — each authority submits their share;
    PVSS verification runs immediately; coalition weight tracked in real time.
  \item \textbf{Tally} — once both reconstruction gates are satisfied,
    the private key is assembled via Lagrange interpolation and all
    ciphertexts are decrypted.
  \item \textbf{Reputation Update} — behaviour scores computed for every
    authority; reputation registry updated; next election's quorum
    pre-computed from new scores.
\end{enumerate}

If reconstruction becomes \emph{impossible} (too many cheaters and skippers
to ever satisfy either gate), the system transitions to a \textbf{failed}
state: automatic penalties are applied, the blockchain records the failure
cause and responsible parties, and the next election is unlocked with an
adjusted (higher) quorum.

% =============================================================================
\section{Reputation Model}
% =============================================================================

\subsection{Behavioural Scores}

After each election, three behavioural dimensions are scored for every
authority $A_i$:

\begin{description}
  \item[Participation ($P_i$)] Did the authority engage with the protocol?
    \begin{equation}
      P_i = \begin{cases} 1.0 & \text{submitted valid share} \\
                          0.5 & \text{submitted invalid share (cheated)} \\
                          0.0 & \text{absent / skipped} \end{cases}
    \end{equation}

  \item[Contribution ($C_i$)] How much did their weight cover the
    per-share average needed?
    \begin{equation}
      C_i = \begin{cases}
        \min\!\left(1,\; \dfrac{w_i}{\bar{w}_{\text{needed}}}\right)
          & \text{if submitted valid share} \\
        0.0 & \text{otherwise}
      \end{cases}
    \end{equation}
    where $\bar{w}_{\text{needed}} = \tau / t$ and $\tau$ is the quorum,
    $t$ the Shamir threshold.

  \item[Honesty ($H_i$)] Did the authority submit the correct share?
    \begin{equation}
      H_i = \begin{cases} 1.0 & \text{PVSS verification passed} \\
                          0.0 & \text{PVSS verification failed (cheated)} \\
                          0.5 & \text{absent — neutral, did not lie} \end{cases}
    \end{equation}
\end{description}

\subsection{Reputation Update Formula}

The reputation $R_i$ is updated after each election using an
exponentially weighted formula:
\begin{equation}
  \boxed{R_i^{(\ell+1)} = \delta \left[ \lambda R_i^{(\ell)} +
    (1-\lambda)\left(\alpha P_i + \beta C_i + \gamma H_i\right) \right]}
  \label{eq:rep}
\end{equation}

\noindent with parameters:
\begin{center}
\begin{tabular}{cll}
\toprule
Parameter & Value & Role \\
\midrule
$\alpha$ & 0.2 & Weight of Participation in new score \\
$\beta$  & 0.3 & Weight of Contribution in new score \\
$\gamma$ & 0.5 & Weight of Honesty in new score \\
$\lambda$ & 0.7 & Memory factor (past vs.\ present) \\
$\delta$ & 0.95 & Decay factor (small penalty for time elapsed) \\
\bottomrule
\end{tabular}
\end{center}

\noindent Honesty carries the highest weight ($\gamma = 0.5$) because
submitting a corrupted share is the most damaging possible action.
Memory ($\lambda = 0.7$) ensures that a single bad election does not
permanently destroy a long-standing reputation, but sustained misbehaviour
accumulates.

\subsection{Initial Reputation}

New authorities begin at:
\begin{equation}
  R_{\text{initial}} = \frac{2}{3}
\end{equation}
This value is chosen so that a group of entirely fresh, equal-reputation
authorities produces a neutral starting quorum of exactly $\lceil N/2 \rceil$
(proven in Section~\ref{sec:threshold}).

\subsection{Reputation Weights}

Scores are normalised into per-authority weights relative to the group:
\begin{equation}
  w_i = \frac{R_i}{\sum_{j=1}^{N} R_j} \cdot N
  \label{eq:weight}
\end{equation}
A neutral group (all equal scores) gives $w_i = 1.0$ for all $i$.
An authority with twice the average reputation has $w_i = 2.0$.
Weights sum to $N$ by construction.

\subsection{Persistence Across Elections}

Reputations are stored in a persistent registry keyed by authority name
(public key in a production deployment). When an authority participates
in a new election, their historical score is loaded automatically.
New authorities who have never participated start at $R_{\text{initial}}$.
This means the system has \emph{memory}: a long track record of honest
behaviour genuinely lowers future quorum requirements.

% =============================================================================
\section{Adaptive Threshold}
\label{sec:threshold}
% =============================================================================

\subsection{Trust and Evenness}

The \emph{collective trust} of a group is measured as the product of
average reputation and a distributional evenness factor:

\begin{equation}
  E = \max\!\left(0,\; 1 - \frac{\sum_i (R_i - \bar{R})^2}{N \bar{R}^2}\right)
  \label{eq:evenness}
\end{equation}

\begin{equation}
  \tau_{\text{trust}} = \min(1,\; \bar{R} \cdot E)
  \label{eq:trust}
\end{equation}

The evenness factor $E$ is close to 1 when all reputations are similar
and drops toward 0 when they are highly unequal. This prevents a scenario
where a single dominant authority with very high reputation (and
correspondingly high weight) could reconstruct alone while the threshold
drops — a concentration attack impossible to execute because:
\begin{property}[Concentration resistance]
  If one authority's weight approaches $N-1$, evenness $E \to 0$, trust
  $\tau_{\text{trust}} \to 0$, and the quorum rises to its maximum $N-1$.
  Therefore no single authority can ever accumulate sufficient weight to
  meet the quorum alone.
\end{property}

\subsection{Two-Piece Linear Adaptive Formula}

The effective quorum $\tau$ is computed via a two-piece linear function
anchored at $\lceil N/2 \rceil$ for neutral trust:

\begin{align}
  \tau_{\text{neutral}} &= \left\lceil \frac{N}{2} \right\rceil \\[4pt]
  \tau_{\text{max}}     &= N - 1 \\[4pt]
  \tau_{\text{min}}     &= \max\!\left(2,\, \left\lceil \frac{N}{3} \right\rceil\right)
\end{align}

\begin{equation}
  \boxed{\tau =
    \begin{cases}
      \tau_{\text{max}} - (\tau_{\text{max}} - \tau_{\text{neutral}})
        \cdot \tau_{\text{trust}} \cdot \tfrac{3}{2}
        & \text{if } \tau_{\text{trust}} \leq \tfrac{2}{3} \\[6pt]
      \tau_{\text{neutral}} - (\tau_{\text{neutral}} - \tau_{\text{min}})
        \cdot (\tau_{\text{trust}} - \tfrac{2}{3}) \cdot 3
        & \text{if } \tau_{\text{trust}} > \tfrac{2}{3}
    \end{cases}}
  \label{eq:threshold}
\end{equation}

\begin{theorem}[Neutral anchor]
  For $N$ authorities all at $R_i = \tfrac{2}{3}$ (the initial reputation),
  Equation~\eqref{eq:threshold} yields $\tau = \lceil N/2 \rceil$ exactly.
\end{theorem}
\begin{proof}
  All equal reputations give $E = 1$ (zero variance), so
  $\tau_{\text{trust}} = \bar{R} \cdot 1 = \tfrac{2}{3}$.
  Substituting into the first branch:
  $\tau = \tau_{\text{max}} - (\tau_{\text{max}} - \tau_{\text{neutral}}) \cdot \tfrac{2}{3} \cdot \tfrac{3}{2}
         = \tau_{\text{max}} - (\tau_{\text{max}} - \tau_{\text{neutral}})
         = \tau_{\text{neutral}} = \lceil N/2 \rceil$.
\end{proof}

\noindent The quorum range for representative group sizes:

\begin{center}
\begin{tabular}{ccccc}
\toprule
$N$ & $\tau_{\text{neutral}}$ & $\tau_{\text{min}}$ (full trust) & $\tau_{\text{max}}$ (no trust) \\
\midrule
3 & 2 & 2 & 2 \\
4 & 2 & 2 & 3 \\
5 & 3 & 2 & 4 \\
6 & 3 & 2 & 5 \\
7 & 4 & 3 & 6 \\
8 & 4 & 3 & 7 \\
\bottomrule
\end{tabular}
\end{center}

% =============================================================================
\section{Two-Gate Reconstruction Policy}
% =============================================================================

Reconstruction is triggered only when \emph{both} conditions are satisfied
simultaneously:

\begin{equation}
  \underbrace{|\mathcal{S}| \geq t}_{\text{Gate 1: cryptographic}}
  \quad \text{AND} \quad
  \underbrace{\sum_{i \in \mathcal{S}} w_i \geq \tau}_{\text{Gate 2: trust}}
  \label{eq:twogates}
\end{equation}

where $\mathcal{S}$ is the set of authorities who submitted valid shares,
$t$ is the Shamir threshold (derived as $\lfloor \tau \rfloor$, minimum 2),
and $\tau$ is the current adaptive quorum.

\subsection{Why Two Gates Are Necessary}

\begin{description}
  \item[Gate 1 alone (share count only)] An adversary controlling many
    low-reputation authorities could submit enough shares to satisfy $t$
    while their combined weight is far below $\tau$. The weight gate blocks
    this \emph{flooding attack}.

  \item[Gate 2 alone (weight only)] A single authority with very high
    reputation might accumulate weight $\geq \tau$ if the threshold dropped
    sufficiently. The share count gate (minimum 2) blocks this
    \emph{monopoly attack}. The evenness mechanism additionally raises $\tau$
    when one authority dominates, providing a second line of defence.
\end{description}

\subsection{Impossibility Detection}

The system continuously monitors whether reconstruction is still achievable.
Let $\mathcal{V}$ be the set of valid-share submitters, $\mathcal{F}$ the
set of authorities who have acted (valid, invalid, or skipped), and
$\mathcal{U} = \{A_i\} \setminus \mathcal{F}$ the remaining free authorities.

\begin{equation}
  \text{Impossible} \iff |\mathcal{V}| + |\mathcal{U}| < t
  \;\;\text{OR}\;\;
  \sum_{i \in \mathcal{V}} w_i + \sum_{i \in \mathcal{U}} w_i < \tau
\end{equation}

If impossible, the system transitions to a \texttt{failed} state
immediately — before waiting for remaining authorities to act —
applies automatic penalties to all responsible parties, and records
the event on the blockchain.

% =============================================================================
\section{Blockchain Integration}
% =============================================================================

\subsection{Role of the Blockchain}

The blockchain in S$^3$-Vote serves as a \textbf{tamper-evident public
bulletin board}. It does not provide consensus (single-node prototype)
but provides the structural guarantees needed for the protocol:

\begin{enumerate}
  \item \textbf{Commitment anchoring} — PVSS commitments are published
    before any share submission, preventing retroactive commitment
    substitution.
  \item \textbf{Vote integrity} — encrypted ballots are written before
    reconstruction begins; the tally cannot be manipulated retroactively.
  \item \textbf{Non-repudiation} — share submissions and failures are
    permanently recorded with the submitting authority's identity;
    a cheater cannot deny the act.
  \item \textbf{Reputation audit trail} — every reputation update is
    recorded with its round number; the full history is independently
    verifiable.
\end{enumerate}

\subsection{Block Structure}

Each block $B_k$ contains:
\begin{itemize}
  \item Index $k$, timestamp
  \item Data payload (event dictionary)
  \item $\text{hash}_k = \text{SHA-256}(\text{index} \,\|\, \text{timestamp}
    \,\|\, \text{data} \,\|\, \text{hash}_{k-1})$
\end{itemize}
Altering any block invalidates all subsequent hashes, making tampering
detectable by any observer who holds a copy of the chain.

\subsection{Events Recorded}

\begin{center}
\begin{tabular}{ll}
\toprule
Event & Payload \\
\midrule
\texttt{election\_start}        & title, options, authorities, quorum, PVSS commitments \\
\texttt{vote\_cast}             & round, ElGamal ciphertext $(C_1, C_2)$ \\
\texttt{voting\_closed}         & round, total votes cast \\
\texttt{share\_submitted}       & round, authority, coalition weight at submission \\
\texttt{share\_invalid}         & round, authority (PVSS verification failed) \\
\texttt{share\_skipped}         & round, authority \\
\texttt{reconstruction\_complete} & round, coalition, coalition weight, tally \\
\texttt{election\_failed}       & round, valid shares, cheaters, skippers \\
\texttt{reputation\_update}     & round, new reputation scores, new quorum \\
\bottomrule
\end{tabular}
\end{center}

\subsection{Production Deployment Note}

The prototype implements the blockchain as an in-memory hash-chained
list on a single node. In a production deployment, this would be
replaced by a distributed ledger (e.g., Ethereum or Hyperledger Fabric),
where no single party controls the chain and tamper-evidence is enforced
by network consensus. The data model and cryptographic properties are
identical; only the storage and consensus layer changes.

% =============================================================================
\section{Security Analysis}
% =============================================================================

\subsection{Vote Confidentiality}

Individual votes remain confidential until reconstruction because:
(a) the ElGamal private key is split across $N$ authorities;
(b) any $t-1$ or fewer shares reveal no information about the key
(information-theoretic security of Shamir's scheme);
(c) ciphertexts on the blockchain are computationally indistinguishable
from random under the DDH assumption.

\subsection{Tally Integrity}

Once voting closes, the \texttt{voting\_closed} block is chained.
Any attempt to add, remove, or modify a \texttt{vote\_cast} block
breaks the hash chain from that point forward, detectable by any
holder of the chain.

\subsection{Cheat Detection and Accountability}

PVSS commitments are written at election creation. A cheating authority
$A_i$ who submits $\hat{s}_i \neq s_i$ will fail the commitment check
$g^{H(\hat{s}_i)} \not\equiv C_i \pmod{p}$. The failure, the authority
identity, and the round are recorded as a \texttt{share\_invalid} block —
public, permanent, and hash-chained into the ledger before $A_i$ can
dispute it.

\subsection{Sybil Resistance}

The reputation system is resistant to Sybil attacks (creating many
fake identities) because new authorities start at $R_{\text{initial}} = 2/3$
and must build a track record over multiple elections to gain elevated
weight. A flood of new identities in a round does not lower the quorum —
it raises it, because many new authorities with neutral reputation
increase $N$ and thus $\tau_{\text{neutral}} = \lceil N/2 \rceil$.

\subsection{Collusion Resistance}

A colluding coalition of $k$ low-reputation authorities cannot satisfy
the weight gate unless $\sum_{i \in \mathcal{C}} w_i \geq \tau$.
Since low-reputation authorities have $w_i < 1$, they need more members
than a trusted coalition — and the evenness factor raises $\tau$ further
when the honest authorities' reputations are heterogeneous.

\subsection{Concentration Attack}

As proven in Property 1, a single authority cannot accumulate enough
weight to reconstruct alone regardless of their reputation, because the
evenness factor $E$ approaches zero as weight concentrates, driving the
quorum to $\tau_{\text{max}} = N-1$.

% =============================================================================
\section{Comparison with Prior Work}
% =============================================================================

\subsection{Against Shamir's Threshold Scheme (1979)}

Shamir's original $(t,n)$-threshold scheme \cite{shamir1979} treats every
share as equal and provides no mechanism for identifying or penalising
misbehaving participants. Key differences:

\begin{center}
\begin{tabular}{lll}
\toprule
Property & Shamir (1979) & S$^3$-Vote \\
\midrule
Share equality & All equal & Reputation-weighted \\
Cheat detection & None & PVSS — immediate, on-chain proof \\
Threshold & Fixed & Adaptive (reputation-driven) \\
Failed round handling & Not addressed & Automatic penalty + recovery \\
Authority accountability & None & Persistent, cross-round reputation \\
\bottomrule
\end{tabular}
\end{center}

\subsection{Against Nojoumian's Social Secret Sharing}

Nojoumian \textit{et al.}\ \cite{nojoumian2012,nojoumian2018} introduced
reputation into secret sharing but left several critical gaps:

\begin{center}
\begin{tabular}{lll}
\toprule
Property & Nojoumian & S$^3$-Vote \\
\midrule
Trusted dealer & Required & None \\
Reputation assignment & Manual / social (subjective) & Formula-based, automatic \\
Cheat detection & Implicit (failure only) & PVSS — public, on-chain \\
Reconstruction gate & Single (weighted shares) & Dual (shares AND weight) \\
Concentration protection & None & Evenness factor \\
Neutral starting quorum & Ad hoc & $\lceil N/2 \rceil$ (proven) \\
Failed election recovery & Not addressed & Automatic penalties + recovery \\
Audit trail & Database & Blockchain (hash-chained) \\
Application domain & General secret sharing & Complete voting system \\
\bottomrule
\end{tabular}
\end{center}

\subsubsection*{Trusted Dealer Elimination}

Nojoumian's scheme requires a trusted dealer to initialise reputation
scores and arbitrate updates. This reintroduces the central trust problem
that threshold cryptography exists to solve.
S$^3$-Vote derives reputation entirely from verifiable on-chain behaviour —
no coordinator, no subjectivity.

\subsubsection*{Objective vs.\ Social Reputation}

In Nojoumian's model, authorities evaluate each other — a social process
that is gameable (coalitions can inflate each other's scores).
In S$^3$-Vote, scores are computed deterministically from Equation
\eqref{eq:rep} applied to PVSS-verified outcomes. There is no way to
inflate a reputation except by genuinely behaving well.

\subsubsection*{Two-Gate Policy}

Nojoumian uses a single adjusted threshold. S$^3$-Vote enforces
Equation \eqref{eq:twogates}: both the cryptographic share count and the
trust weight must be satisfied. This closes the flooding attack (many
low-weight shares) and the monopoly attack (one very high-weight authority)
that single-gate schemes cannot prevent.

\subsubsection*{Provable Neutral Anchor}

Nojoumian does not derive a principled starting threshold. S$^3$-Vote
proves (Theorem 1) that fresh authorities always begin at the majority
quorum $\lceil N/2 \rceil$, providing a principled, consistent baseline
regardless of group size.

% =============================================================================
\section{Prototype Implementation}
% =============================================================================

The S$^3$-Vote prototype is implemented in Python (backend) and
JavaScript (single-page frontend), comprising the following modules:

\begin{description}
  \item[\texttt{crypto/encryption.py}] ElGamal key generation, vote
    encryption and decryption.
  \item[\texttt{crypto/shamir.py}] Polynomial-based $(t,n)$-threshold
    secret splitting and Lagrange interpolation reconstruction.
  \item[\texttt{crypto/pvss.py}] Commitment generation and share
    verification.
  \item[\texttt{reputation/reputation\_engine.py}] Implementation of
    Equation \eqref{eq:rep}.
  \item[\texttt{reputation/threshold.py}] Implementation of Equations
    \eqref{eq:evenness}–\eqref{eq:threshold}.
  \item[\texttt{blockchain/blockchain.py}] Hash-chained block structure
    with SHA-256 linking.
  \item[\texttt{api/app.py}] Flask REST API orchestrating the full
    election lifecycle.
  \item[\texttt{api/templates/index.html}] Web dashboard for election
    management, vote simulation, share submission, and reputation
    visualisation.
\end{description}

The dashboard allows an operator to create elections with arbitrary
authority sets, simulate voting, and control each authority's share
submission (valid, invalid, or skipped) to observe the quorum and
reputation dynamics in real time. Reputation persists across sessions
via the in-memory registry and reflects accurately in the quorum
computation for subsequent elections.

% =============================================================================
\section{Example Walkthrough}
% =============================================================================

We trace a two-round scenario with $N = 7$ authorities.

\subsubsection*{Round 1 — Initial State}

All 7 authorities start at $R_i = 2/3$.
By Theorem 1, the quorum is $\tau = \lceil 7/2 \rceil = 4$.
Suppose 5 submit valid shares, 1 cheats, and 1 skips.
The coalition of 5 valid submitters achieves weight $5 \times 1.0 = 5.0 \geq 4$
and share count $5 \geq t = 4$.
Reconstruction succeeds.

After the round, reputation updates via Equation \eqref{eq:rep}:
\begin{align*}
  R_{\text{honest}} &\approx 0.95[0.7(2/3) + 0.3(1.0)] \approx 0.728 \\
  R_{\text{cheater}} &\approx 0.95[0.7(2/3) + 0.3(0.5 \cdot 0.2 + 0 \cdot 0.3 + 0 \cdot 0.5)] \approx 0.472 \\
  R_{\text{skipper}} &\approx 0.95[0.7(2/3) + 0.3(0 \cdot 0.2 + 0 \cdot 0.3 + 0.5 \cdot 0.5)] \approx 0.515
\end{align*}

\subsubsection*{Round 2 — Reputations Loaded from Registry}

The new reputations $\{0.728 \times 5,\, 0.472,\, 0.515\}$ are loaded.
The honest majority has higher average rep but the two underperformers
reduce evenness.
Trust $\tau_{\text{trust}} \approx 0.67$ remains near neutral, so the
quorum adjusts to approximately $\tau \approx 4.07$ — slightly \emph{above}
neutral, reflecting that the group is less cohesive than before.
The cheater and skipper must now work against a marginally stricter quorum,
while repeated honest behaviour across subsequent rounds will drive it
back down toward 3.

% =============================================================================
\section{Limitations and Future Work}
% =============================================================================

\begin{description}
  \item[Single-node blockchain] The prototype runs on one server.
    Production deployment requires a distributed ledger with consensus.
  \item[Simulated authority clients] All authority actions are simulated
    by one administrator in the prototype. A production system gives each
    authority an independent client that holds their private key and signs
    share submissions.
  \item[Cryptographic identity] Authority names are currently plain strings.
    In production, each name maps to a public key; share submissions are
    signed by the corresponding private key and verified on-chain, preventing
    impersonation.
  \item[Parameter sensitivity] The reputation parameters
    $(\alpha, \beta, \gamma, \lambda, \delta)$ are set by design rationale.
    Empirical calibration across realistic election scenarios would
    strengthen the parameter choices.
  \item[Formal security proof] A rigorous game-based security proof under
    standard cryptographic assumptions (DDH, random oracle) is left for
    future work.
\end{description}

% =============================================================================
\section{Conclusion}
% =============================================================================

S$^3$-Vote demonstrates that threshold-cryptographic voting can be made
self-governing: no trusted dealer, no static parameters, no silent failures.
By anchoring PVSS commitments and all authority actions on a blockchain,
every act of honesty or betrayal becomes a permanent, public, verifiable
record. By feeding that record into a reputation formula and propagating
reputations across elections, the system develops institutional memory.
And by translating that memory into an adaptive two-gate quorum, it turns
good behaviour into a concrete, quantifiable reward — fewer participants
needed to reconstruct — and bad behaviour into a concrete cost — a higher,
stricter quorum for everyone.

The core theoretical contribution is the proof that the neutral quorum
is always $\lceil N/2 \rceil$ regardless of group size, that evenness
prevents weight concentration, and that the dual-gate policy closes
attack vectors that prior single-threshold schemes cannot address.

Compared to Shamir's foundational scheme, S$^3$-Vote adds accountability,
verifiability, and adaptivity.
Compared to Nojoumian's social secret sharing, it eliminates the trusted
dealer, replaces subjective social reputation with objective on-chain
evidence, introduces dual-gate reconstruction, and provides a complete,
deployable voting system rather than a secret sharing primitive.

% =============================================================================
\bibliographystyle{plain}
\begin{thebibliography}{9}

\bibitem{shamir1979}
A.\ Shamir,
``How to Share a Secret,''
\textit{Communications of the ACM}, vol.\ 22, no.\ 11, pp.\ 612--613, 1979.

\bibitem{elgamal1985}
T.\ ElGamal,
``A Public Key Cryptosystem and a Signature Scheme Based on Discrete Logarithms,''
\textit{IEEE Transactions on Information Theory}, vol.\ 31, no.\ 4,
pp.\ 469--472, 1985.

\bibitem{schoenmakers1999}
B.\ Schoenmakers,
``A Simple Publicly Verifiable Secret Sharing Scheme and its Application to
Electronic Voting,''
in \textit{Advances in Cryptology — CRYPTO 1999}, Lecture Notes in Computer
Science, vol.\ 1666, pp.\ 148--164, Springer, 1999.

\bibitem{nojoumian2012}
M.\ Nojoumian and D.\ R.\ Stinson,
``Social Secret Sharing in Cloud Computing Using a New Trust Function,''
in \textit{Proceedings of the 10th Annual Conference on Privacy, Security
and Trust (PST)}, pp.\ 161--167, IEEE, 2012.

\bibitem{nojoumian2018}
M.\ Nojoumian,
``Novel Secret Sharing and Commitment Schemes for Cryptographic Applications,''
Ph.D.\ dissertation, University of Waterloo, 2012.
(Extended work published through 2018.)

\bibitem{pedersen1991}
T.\ P.\ Pedersen,
``Non-Interactive and Information-Theoretic Secure Verifiable Secret Sharing,''
in \textit{Advances in Cryptology — CRYPTO 1991}, Lecture Notes in Computer
Science, vol.\ 576, pp.\ 129--140, Springer, 1991.

\bibitem{desmedt1994}
Y.\ Desmedt and Y.\ Frankel,
``Threshold Cryptosystems,''
in \textit{Advances in Cryptology — CRYPTO 1989}, Lecture Notes in Computer
Science, vol.\ 435, pp.\ 307--315, Springer, 1990.

\bibitem{cramer1997}
R.\ Cramer, R.\ Gennaro, and B.\ Schoenmakers,
``A Secure and Optimally Efficient Multi-Authority Election Scheme,''
in \textit{Advances in Cryptology — EUROCRYPT 1997}, Lecture Notes in
Computer Science, vol.\ 1233, pp.\ 103--118, Springer, 1997.

\bibitem{chaum1988}
D.\ Chaum, C.\ Cr{\'{e}}peau, and I.\ Damg{\aa}rd,
``Multiparty Unconditionally Secure Protocols,''
in \textit{Proceedings of the 20th Annual ACM Symposium on Theory of
Computing (STOC)}, pp.\ 11--19, ACM, 1988.

\end{thebibliography}

\end{document}
